\begin{table}[t]
\centering
\begin{threeparttable}
\caption{Inference procedures for the baseline estimate under a single treated unit}
\label{tab:inference}
\begin{tabular}{lc}
\toprule
 & Baseline specification \\
\midrule
$\hat\tau$: BTC $\times$ Post (Fisher-$z$) & -0.0212 \\
 & (0.0170) \\
$\hat\tau$ translated into correlation units & -0.0181 \\
Bitcoin pre-period mean correlation & 0.3822 \\
Randomization $p$ (space $\times$ time grid) --- primary & 0.665 \\
Randomization $p$ (placebo-in-space only) & 0.80 \\
Wild cluster bootstrap $p$ (unreliable here) & 0.467 \\
Placebo assignments in the grid & 199 \\
SD of the randomization distribution & 0.0619 \\
Randomization 95\% CI, lower (Fisher-$z$) & -0.1426 \\
Randomization 95\% CI, upper (Fisher-$z$) & 0.1002 \\
Conley--Taber 95\% CI, lower (Fisher-$z$) & -0.1065 \\
Conley--Taber 95\% CI, upper (Fisher-$z$) & 0.0386 \\
$N$ & 540 \\
\bottomrule
\end{tabular}
\begin{tablenotes}[flushleft]
\footnotesize
\item All rows refer to the baseline coin-month specification reported in column (1) of Table~\ref{tab:main}. With ten clusters of which one is treated, cluster-robust variance estimation has no valid asymptotic justification and the wild cluster bootstrap is also unreliable, so the randomization and Conley--Taber procedures are the ones on which the paper's claims rest. The randomization $p$-value is the share of placebo assignments producing an absolute estimate at least as large as the actual one, over a grid combining nine placebo coins at the true date with placebo dates drawn from every month of 2021m7--2023m1 on the pre-period sample. The cluster-robust interval printed in Table~\ref{tab:main} is roughly a third the width of the two valid intervals shown here, which is a direct measurement of the downward bias a single treated cluster induces. Correlation units are obtained via $\Delta\rho \approx (1-\bar\rho_{\mathrm{pre}}^2)\,\hat\tau$.
\end{tablenotes}
\end{threeparttable}
\end{table}
